% rexx-in-latex.tex -- ooRexx inside LuaLaTeX, through the Lua module
% (require "oorexx").  How to run it: README.md, next to this file:
%   lualatex --shell-escape rexx-in-latex.tex
% with oorexx.so (oorexx.dll) and lua.cls where LuaTeX's Lua finds them.
\documentclass{article}
\usepackage{booktabs}

% The module, once.  Then \rexx{expression}: a Rexx expression, evaluated by
% rexx.eval; its value goes into the page.  \detokenize keeps the expression
% as it is written (a ~ stays a ~, not LaTeX's tie).
\directlua{rexx = require "oorexx"}
\newcommand\rexx[1]{\directlua{tex.sprint(rexx.eval("\luaescapestring{\detokenize{#1}}"))}}

\begin{document}

\section*{ooRexx in Lua\LaTeX}

\paragraph{Inline.} Rexx's decimal arithmetic in the text: with Rexx's
default nine digits, $2^{100}$ is \rexx{2**100} and $1/7$ is \rexx{1/7}. A
price keeps its zero: \rexx{format(3.1, , 2)} euros. And the word ``bridge''
backwards is \rexx{'bridge'~reverse}.

\paragraph{A program.} \texttt{powers.rex} writes the table below into this
document: each of its lines of \TeX\ goes out through \texttt{ADDRESS LUA},
as Lua code that calls \texttt{tex.sprint} in Lua\TeX's own Lua. With
\texttt{NUMERIC DIGITS 40}, every value is exact.

\begin{center}
\directlua{luatexVersion = rexx.run("powers.rex", 40)}
\end{center}

\paragraph{Both ways.} \texttt{powers.rex} also asked Lua\TeX\ for its
version (\texttt{ADDRESS LUA} returns what the Lua code returns, in
\texttt{RESULT}) and gave it back to this document: Lua\TeX\
\directlua{tex.sprint(luatexVersion)}.

\end{document}
